\documentclass[a4paper]{article}
\usepackage[utf8]{inputenc}
\usepackage{amsmath}
\usepackage{graphicx}
\usepackage[bottom=2.0cm,top=2.0cm,left=2.0cm,right=2.0cm]{geometry}


\title{%
\begin{figure}[htp]
    \centering
    \includegraphics[width=5cm]{logo_unb.png}
\end{figure}
  Métodos Matemáticos da Física \\
  \large Departamento de Matemática - Universidade de Brasília}

\author{
    Felipe Martins Goulart - 15/0153651 \\
    Gabriel Adiron Ribeiro - 17/0103196 \\
}
\date{Março 2021}

\begin{document}
    \maketitle
    \renewcommand{\theenumi}{\alph{}}
    \section {É Verdadeiro(V) ou Falso(F)?}
     \begin{enumerate}
       \begin{enumerate}
         \item Desenvolvendo a  função $ f(x) = x^2, - \pi  < x < \pi$, em série de Fourier com período de $2\pi$ e utilizando o Teorema de Fourier tem-se que
            \begin{equation}
                \begin{aligned}
                    {\frac{1}{1^2}+\frac{1}{2^2}+...+\frac{1}{n^2}+...=\frac{\pi^2}{6}.}\\
                \end{aligned}
            \end{equation}
         \item Existe função $f \epsilon \mathit{L}^2([-\pi,\pi])$ de período $2\pi$ de tal modo que sua ´serie de Fourier seja $\sum_{n=1}^{- \infty} n^2sin(nx)$.
       \end{enumerate}
     \end{enumerate}
    \section{Demonstre que \{$\sqrt{\frac{2}{L}} sin\frac{k\pi}{L}x \}_{k \epsilon N}$ é um \textbf{C.O.N.C.} em $C_{per} ([0,L])$.}
    \section{Demonstre que se $f \epsilon \mathit{L}^2([-L,L])$ então $\frac{a_0^2}{2} +\sum_{n=1}^{- \infty} \a_k^2 + b_k^2 = \frac{1}{L} \int_{-L}^{L} |f(x)|^2dx$, onde $a_0$, $a_k$, $b_{k}$ são coeficientes de Fourier.} 
    \section{Encontre a série de Fourier da função f(x) = cos^5x }
    \section{Explique porque a série $\sum_{n=1}^{- \infty} \frac{cos(nx)}{n}$ não é a série de Fourier de uma função diferneciável}
\end{document}
